\documentclass[11pt]{article}
\usepackage[utf8]{inputenc}
\usepackage{geometry}
\geometry{a4paper, margin=1in}
\usepackage{graphicx}
\usepackage{amsmath}
\usepackage{amssymb}
\usepackage{amsfonts}
\usepackage{booktabs}
\usepackage{float}
\usepackage{hyperref}
\usepackage{listings}
\usepackage{xcolor}
\usepackage{caption}
\usepackage{subcaption}
\usepackage{enumitem}
\usepackage{appendix}

% Define colors for code listings
\definecolor{codegreen}{rgb}{0,0.6,0}
\definecolor{codegray}{rgb}{0.5,0.5,0.5}
\definecolor{codepurple}{rgb}{0.58,0,0.82}
\definecolor{backcolour}{rgb}{0.95,0.95,0.92}

\lstdefinestyle{mystyle}{
    backgroundcolor=\color{backcolour},
    commentstyle=\color{codegreen},
    keywordstyle=\color{magenta},
    numberstyle=\tiny\color{codegray},
    stringstyle=\color{codepurple},
    basicstyle=\ttfamily\footnotesize,
    breakatwhitespace=false,
    breaklines=true,
    captionpos=b,
    keepspaces=true,
    numbers=left,
    numbersep=5pt,
    showspaces=false,
    showstringspaces=false,
    showtabs=false,
    tabsize=2
}
\lstset{style=mystyle}
\title{\textbf{Project 1: Co-creativity as Curation} \\
\large MovieLens 100K Ranking \& Personalization Pipeline}

\author{
Tiago Silva\thanks{\texttt{tds@student.dei.uc.pt}}
\and
Miguel Cabral Pinto\thanks{\texttt{mcp@student.uc.pt}}
\and
Gonçalo Borges\thanks{\texttt{goncaloborges@student.dei.uc.pt}}
\and
André Cardoso\thanks{\texttt{}}
\and
Miguel Martins\thanks{\texttt{}}
}

\date{\today}

\begin{document}

\maketitle

\begin{abstract}

\end{abstract}

\section{Introduction}
\label{sec:introduction}

In modern entertainment platforms, the dominant creative act is often not the generation of a single artifact, but rather the curation, sequencing, and contextualization of what users experience next. This project explores this paradigm by framing curation as a form of co-creation. When a system recommends content, it applies selection pressure that ultimately shapes user tastes, creator incentives, and the visibility of different cultural niches. 

Historically, recommender systems focused heavily on optimizing for prediction accuracy, such as minimizing the Root Mean Square Error (RMSE) of user ratings. However, optimizing solely for accuracy often leads to degenerate system behaviors, such as filter bubbles, popularity bias, and a lack of discovery. A truly co-creative pipeline must balance immediate relevance with list-level properties like diversity and pacing, while continuously adapting to user feedback.

To demonstrate this, we implemented an end-to-end interactive curation pipeline using the MovieLens 100K dataset. Our system blueprint is divided into four distinct stages:

\begin{enumerate}
    \item \textbf{Candidate Generation:} Filtering the dataset to produce a viable pool of unseen movies for a given user.
    \item \textbf{Ranking:} Scoring candidates to encode personal value and relevance. We compare four rankers: a Popularity baseline, Matrix Factorization trained via Stochastic Gradient Descent (MF-SGD), Matrix Factorization trained via Alternating Least Squares (MF-ALS), and a Pairwise Learning-to-Rank (LTR) model.
    \item \textbf{Reranking:} Applying Maximal Marginal Relevance (MMR) to the top-scored items to enforce diversity constraints, ensuring the final recommended slate explores different genres rather than collapsing into a narrow niche.
    \item \textbf{Personalization Update:} Simulating interactive user sessions where the user's latent taste profile is dynamically updated after each choice using an Exponential Moving Average (EMA).
\end{enumerate}

Through this pipeline, this report investigates the comparative offline performance of these rankers (Recall@10 and NDCG@10), analyzes the critical tradeoff between relevance and diversity, and observes how a user's state evolves during simulated co-creative sessions.


\section{Dataset and Preprocessing}
\label{sec:dataset}

For this project, we utilized the \textbf{MovieLens 100K} dataset, a standard benchmark in recommender systems research that contains 100,000 ratings (ranging from 1 to 5) applied to 1,682 movies by 943 users.

To prepare the data for our models, we performed the following preprocessing steps:
\begin{enumerate}
    \item \textbf{ID Normalization:} The original \texttt{user\_id} and \texttt{item\_id} values were mapped to contiguous integer indices starting from 0 (e.g., $u \in \{0, \dots, 942\}$ and $i \in \{0, \dots, 1681\}$). This allows us to use the indices directly for fast array lookups and matrix operations during training.
    \item \textbf{Interaction Histories:} We built per-user and per-item interaction dictionaries. We also calculated the global mean rating ($\mu$) across the training set, as well as the absolute popularity of each item (interaction count), which serves as a baseline feature for our rankers.
    \item \textbf{Data Splitting:} We applied a randomized \textbf{80/20 train/test split} over the observed ratings. Furthermore, to facilitate rigorous hyperparameter tuning without leaking test data, we applied a secondary 80/20 split on the training data to generate a dedicated \textbf{validation set}.
    \item \textbf{Relevance Definition:} For the purposes of evaluating our Top-K recommendation lists, we defined the set of "relevant" test items for a user $u$ mathematically as:
    $$G_u = \{i : (u,i) \in \Omega_{test} \text{ and } r_{ui} \ge 4\}$$
    where $\Omega_{test}$ represents the test set and $r_{ui}$ is the true rating given by the user.
\end{enumerate}

During evaluation, to ensure we are strictly evaluating the system's ability to recommend \textit{unseen} content, all items that a user had already rated in the training set were excluded from their candidate generation pool.

\section{Implementation Details}
\label{sec:implementation}

The co-creative system blueprint is structured into four distinct stages: candidate generation, ranking, applying list constraints (reranking), and a personalization update via user feedback.

\subsection{Candidate Generation}
For any given user $u$, the candidate set $C(u)$ consists of all items that the user has not interacted with in the training split. This ensures that the system only recommends unseen movies to the user.

\subsection{Rankers}
We implemented four distinct ranking methods to encode personal relevance:
\begin{itemize}
    \item \textbf{Popularity Baseline:} A non-personalized ranker that orders candidate movies based on their global interaction count in the training dataset.
    \item \textbf{MF-SGD:} A Matrix Factorization model with biases predicting ratings as $\hat{r}_{ui} = \mu + b_u + b_i + p_u^\top q_i$, trained using Stochastic Gradient Descent.
    \item \textbf{MF-ALS:} A Matrix Factorization model with biases trained using Alternating Least Squares, which relies on closed-form ridge regression updates.
    \item \textbf{Pairwise LTR:} A logistic pairwise preference model leveraging gradient descent. The ranking score is computed as $s_\phi(u,i) = \phi^\top f(u,i)$. The feature vector $f(u,i)$ includes the base MF score, user bias, item bias, the log-normalized popularity of the item, and a constant bias term.
\end{itemize}

\subsection{Hyperparameter Tuning}
While standard baseline parameters (e.g., $d=20$, $\lambda=0.05$) were initially considered, we utilized the \texttt{Optuna} framework to rigorously optimize our Matrix Factorization models against our dedicated validation set. 

Rather than relying on manual guesses, we ran automated study trials to minimize the validation Root Mean Square Error (RMSE). For example, the optimization for the MF-SGD model yielded an ideal latent dimension of $d=81$, factor regularization $\lambda_{factors} \approx 0.1138$, bias regularization $\lambda_{biases} \approx 0.0005$, a learning rate $\eta \approx 0.0116$, and $75$ training epochs. These mathematically tuned parameters were subsequently locked in and used for all final offline evaluations and co-creative simulations.

\subsection{Reranking (Diversity Constraints)}
To control list-level diversity, we applied Maximal Marginal Relevance (MMR) over the top 80 candidates produced by the base ranker. The reranker dynamically builds the slate step-by-step by selecting the item $c^*$ that maximizes:
$$c^* = \arg\max_{c \in C \setminus S} \left( (1 - \alpha)\text{sim}(c, q) - \alpha \max_{s \in S}\text{sim}(c, s) \right)$$
Our pipeline computes item similarity via the cosine similarity of the $L_2$-normalized genre vectors extracted from the dataset's item metadata. We explored the trade-off parameter $\alpha \in \{0.1, 0.4, 0.7\}$ to observe the impact on relevance and diversity metrics.
\subsection{Personalization Update (Session)}
To simulate a co-creative loop, the system executes multi-turn sessions where a user is presented with a slate, makes a choice, and has their latent state updated. We simulate the choice by selecting an item from the slate that the user actually rated in the test set. The personalization update applies an Exponential Moving Average (EMA) to the user's latent representation. If user $u$ selects item $i^*$, their factor vector $p_u$ is updated dynamically using the learned item factor $q_{i^*}$ with a rate $\beta = 0.1$:
$$p_u \leftarrow (1 - \beta)p_u + \beta q_{i^*}$$
In our demonstration, we simulated 5 consecutive feedback rounds per user for 3 distinct users to observe how the recommended slate evolves and adapts to recent choices.

\section{Experiments and Results}
\label{sec:experiments}

\subsection{E1: Baseline Comparison}
\label{sec:e1}



\section{Conclusion}
\label{sec:conclusion}

[Summarize the main findings of your project. What did you learn about the tradeoffs between relevance and diversity? Which ranking method performed best? Reflect on the challenges of offline evaluation and the importance of co-creative systems.]

\begin{appendices}
\section{Theory Questions}
\label{app:theory}

\subsection*{1. Stochastic Gradient Descent (SGD)}

We consider the biased matrix factorization model

\begin{align*}
    \mathcal{E}_{ui} &= R_{ui} - \hat{R}_{ui}
    \quad \text{(prediction error)} \\
    \hat{R}_{ui} &= \mu + b_u + b_i + p_u^T q_i
    \quad \text{(prediction model)} \\
    l_{ui} &= (R_{ui} - \hat{R}_{ui})^2
    + \lambda (||p_u||^2 + ||q_i||^2 + b_u^2 + b_i^2)
    \quad \text{(regularized squared loss)}
\end{align*}

The SGD update rule for each parameter $\theta$ is

\begin{align*}
    \theta \leftarrow \theta - \eta \nabla_{\theta}l_{ui},
\end{align*}

where $\eta$ denotes the learning rate.

Using the chain rule,

\begin{align*}
    \frac{\partial}{\partial\theta}(R_{ui} - \hat{R}_{ui})^2
    &= 2(R_{ui} - \hat{R}_{ui})
    \cdot
    \frac{\partial}{\partial\theta}(R_{ui} - \hat{R}_{ui}) \
    &= -2 \mathcal{E}_{ui}
    \cdot
    \frac{\partial\hat{R}_{ui}}{\partial\theta}.
\end{align*}

\paragraph{Gradient with respect to $p_u$}

Only the inner-product term depends on $p_u$, therefore

\begin{align*}
    \nabla_{p_u} l_{ui}
    &= -2\mathcal{E}_{ui} q_i + 2\lambda p_u.
\end{align*}

Thus the SGD update becomes

\begin{align*}
    p_u \leftarrow p_u + \eta(\mathcal{E}_{ui} q_i - \lambda p_u).
\end{align*}

\paragraph{Gradient with respect to $q_i$}

Similarly,

\begin{align*}
    \nabla_{q_i} l_{ui}
    &= -2\mathcal{E}_{ui} p_u + 2\lambda q_i,
\end{align*}

leading to

\begin{align*}
    q_i \leftarrow q_i + \eta(\mathcal{E}_{ui} p_u - \lambda q_i).
\end{align*}

\paragraph{Gradient with respect to $b_u$}

Since the prediction depends linearly on $b_u$,

\begin{align*}
    \nabla_{b_u} l_{ui}
    &= -2\mathcal{E}_{ui} + 2\lambda b_u,
\end{align*}

and therefore

\begin{align*}
    b_u \leftarrow b_u + \eta(\mathcal{E}_{ui} - \lambda b_u).
\end{align*}

\paragraph{Gradient with respect to $b_i$}

Analogously,

\begin{align*}
    \nabla_{b_i} l_{ui}
    &= -2\mathcal{E}_{ui} + 2\lambda b_i,
\end{align*}

so that

\begin{align*}
    b_i \leftarrow b_i + \eta(\mathcal{E}_{ui} - \lambda b_i).
\end{align*}

These updates define the SGD optimization procedure applied independently to each observed rating pair $(u,i)$.

\vspace{1cm}

\subsection*{2. Alternating Least Squares (ALS)}

Assume $Q$, $b_u$, $b_i$, and $\mu$ are fixed. Define the bias-corrected residual

\begin{align*}
    e_{ui} = R_{ui} - (\mu + b_u + b_i).
\end{align*}

Then optimizing over $p_u$ reduces to

\begin{align*}
    L(p_u)
    =
    \sum_{i \in \Omega(u)}
    (e_{ui} - p_u^T q_i)^2
    +
    \lambda ||p_u||^2,
\end{align*}

which corresponds to a ridge regression problem with predictors $q_i$ and targets $e_{ui}$.

Taking the gradient with respect to $p_u$,

\begin{align*}
    \nabla_{p_u} L(p_u)
    &=
    \sum_{i \in \Omega(u)}
    2(e_{ui} - p_u^T q_i)(-q_i)
    +
    2\lambda p_u.
\end{align*}

Setting the gradient equal to zero yields the normal equations

\begin{align*}
    \sum_{i \in \Omega(u)}
    (e_{ui} - p_u^T q_i)(-q_i)
    + \lambda p_u = 0.
\end{align*}

Expanding terms,

\begin{align*}
    \sum_{i \in \Omega(u)}
    e_{ui} q_i
    =
    \sum_{i \in \Omega(u)}
    (q_i q_i^T)p_u
    +
    \lambda p_u.
\end{align*}

Factoring $p_u$,

\begin{align*}
    \sum_{i \in \Omega(u)}
    e_{ui} q_i
    =
    \left(
    \sum_{i \in \Omega(u)}
    q_i q_i^T
    +
    \lambda I
    \right)
    p_u.
\end{align*}

Multiplying both sides by the inverse matrix gives the closed-form solution

\begin{align*}
    p_u
    =
    \left(
    \sum_{i \in \Omega(u)}
    q_i q_i^T
    +
    \lambda I
    \right)^{-1}
    \left(
    \sum_{i \in \Omega(u)}
    e_{ui} q_i
    \right).
\end{align*}

This expression corresponds to the ALS update step for user latent factors under $\ell_2$ regularization.

\vspace{1cm}

\subsection*{3. Pairwise gradient}

Let

\begin{align*}
    s_\phi(u,i) &= \phi^T f(u,i), \\
    \Delta f &= f(u,i^+) - f(u,i^-).
\end{align*}

We want to compute

\begin{align*}
    \nabla_\phi \left(-\log \sigma(\phi^T \Delta f)\right).
\end{align*}

Define

\begin{align*}
    z = \phi^T \Delta f.
\end{align*}

Since

\begin{align*}
    \sigma(z) = \frac{1}{1+e^{-z}},
\end{align*}

its derivative is

\begin{align*}
    \sigma'(z) = \sigma(z)(1-\sigma(z)).
\end{align*}

Therefore,

\begin{align*}
    \frac{d}{dz}\left(-\log\sigma(z)\right)
    &=
    -\frac{1}{\sigma(z)}\sigma'(z) \\
    &=
    -(1-\sigma(z)) \\
    &=
    \sigma(z)-1.
\end{align*}

Using the chain rule,

\begin{align*}
    \nabla_\phi z = \nabla_\phi (\phi^T \Delta f) = \Delta f.
\end{align*}

Thus,

\begin{align*}
    \nabla_\phi \left(-\log\sigma(\phi^T \Delta f)\right)
    =
    (\sigma(\phi^T \Delta f)-1)\Delta f.
\end{align*}


\vspace{1cm}

\subsection*{4. Why RMSE $\neq$ ranking quality?}

Optimizing RMSE in matrix factorization minimizes prediction error over rating values:

\begin{align*}
\text{RMSE}
    =
    \sqrt{
    \frac{1}{|\Omega|}
    \sum_{(u,i)\in\Omega}
    (R_{ui}-\hat{R}_{ui})^2
    }.
\end{align*}

However, top-$K$ recommendation quality depends primarily on the correctness of the ordering of predicted scores rather than their absolute values.

There are several reasons for this mismatch:

\begin{itemize}

\item RMSE penalizes magnitude errors instead of ranking errors. Small prediction changes can significantly affect ranking while barely changing RMSE.

\item It treats all prediction errors equally, whereas recommendation quality depends mainly on correctly ranking the highest-scored items.

\item The absence of a rating is typically related to its value. Users rate only items they are exposed to, introducing exposure bias into the dataset.

\item It does not enforce pairwise ordering constraints such as

\begin{align*}
\hat{R}_{ui^+} > \hat{R}_{ui^-},
\end{align*}

which is essential for ranking-based recommendation tasks.

\end{itemize}
\begin{align*}
\text{Thus, in general, low RMSE does not necessarily imply good top-}K\text{ recommendation quality.}
\end{align*}


\vspace{1cm}

\subsection*{5. Offline evaluation caveat}

Offline evaluation assumes that observed interactions reflect true user preferences. However, this assumption fails when exposure is not random.

We can define interaction as,

\begin{align*}
\text{interaction} = \text{exposure} + \text{preference}.
\end{align*}

Users can only interact with items they are exposed to. Therefore, interaction data reflects both user preferences and historical exposure policies.

As a consequence:

\begin{itemize}

\item Popular items are overrepresented in the dataset, leading to popularity bias.

\item New or rarely exposed items appear artificially worse during evaluation.

\item Offline evaluation lacks counterfactual information about how users would behave if exposed to different items.

\end{itemize}

Thus, offline evaluation can produce misleading estimates of recommendation quality unless exposure is randomized or corrected using debiasing techniques, such as inverse propensity scoring, in which data is reweighted to account for the probability of exposure.

\end{appendices}

\end{document}

